\section{Problem formulation}\label{sec:problem}

An \emph{entity} $e$ is a price series identified by
(state, district, market, commodity, variety).
Grade is excluded from the key for the reason given in Section~\ref{sec:data}.
Entity $e$ publishes modal prices $P_{e,t}$ (INR per quintal) on irregular report dates $t$.
For an origin $t$ at which $e$ reports, and a horizon $h\in\{3,7,14\}$ days, the target is the first report of $e$ dated in $[t+h,\,t+h+3]$.
If no report falls in that window, the origin has no label for $h$.
The \emph{next} horizon targets the next report, provided it is at most 7 days away.
Models predict the log return $y_{e,t,h}=\ln P_{e,\tau}-\ln P_{e,t}$, where $\tau$ is the target date.
Forecasts in rupees are $\hat P=P_{e,t}\exp(\hat y)$, and every error metric is computed on that scale.
Using a log-return target lets one global model serve commodities whose price levels differ by two orders of magnitude.

Three tasks are kept separate:
\begin{itemize}
  \item \textbf{forecasting}, the subject of this paper's main results;
  \item \textbf{anomaly detection}, which flags reports that deviate from the entity's trailing behaviour;
  \item \textbf{decision support}, a transparent rule on top of forecast intervals.
\end{itemize}
None of these models causal effects.

\section{Data and provenance}\label{sec:data}

\paragraph{Source.}
The primary source is an archive of AGMARKNET daily reports for 2001-01-10 to 2026-04-21, distributed as a Kaggle dataset
(\texttt{khandelwalmanas/daily-commodity-prices-india}, version~134, published under GODL-India).
Our local copy matches that release file by file in byte size; content hashes were not compared.
Its columns match the official data.gov.in resource \emph{Variety-wise Daily Market Prices Data of Commodity}.
The live data.gov.in API (resource \texttt{9ef84268-\ldots}) needs a key, which we did not have; its adapter is implemented and tested offline.
The archive has \num{75984017} rows.
After quality classification and business-key de-duplication, the silver layer has \num{75644801} rows; 97.6\% of these are internally consistent (\texttt{VALID}).

\paragraph{Scope.}
We use 36 staple commodities: vegetables, fruits, cereals, pulses, oilseeds and cotton (\num{33116234} valid reports since October 2014).
They were chosen from training-period data with four steps:
\begin{enumerate}
  \item take an a-priori list of staple crop-name patterns;
  \item keep matches with more than \num{50000} valid reports in 2021--2023;
  \item drop two duplicate or processed forms;
  \item keep the 36 most-reported.
\end{enumerate}
This is not the overall top 36: six high-volume items outside the pattern list, such as radish, papaya and jaggery, are excluded.

\paragraph{Identifier harmonisation.}
Two changes in how the source publishes identifiers would otherwise break series:
\begin{itemize}
  \item \textbf{Grade relabel.} From 2025-01-27 many series switch their grade label from \emph{FAQ} to \emph{Local}, with no overlap.
    Across all (state, district, market, commodity, variety, date) keys since 2023, none carries two grades on the same day, so we treat grade as an attribute.
  \item \textbf{Market rename.} With the AGMARKNET~2.0 migration (portal live 2025-11-07), 95--97\% of reports in December 2025 to March 2026 (77\% in April) carry an `` APMC'' suffix on the market name, which we strip.
\end{itemize}
Harmonisation merges \num{117623} fragmented series into \num{75814} entities.
It removes \num{191} same-day collisions.
Without it, essentially no panel series would continue into 2026.
Around the migration, monthly volume collapses to about 6\% of its earlier level in November 2025 and recovers only to roughly 20--35\% in December 2025 to April 2026, so we report test metrics separately before and after November 2025.

\paragraph{Cleaning.}
A causal sanity rule removes reports that differ by more than a factor of five from the geometric mean of the entity's previous ten reports.
These are typically unit or keying slips.
The rule removes \num{28567} reports (0.086\%).
Genuine price spikes below that factor are kept.

\paragraph{Panel.}
An entity enters the training and evaluation panel if it has at least 300 report days in 2021--2023.
This rule uses training-period information only.
The panel has \num{11845} entities and \num{21.06}M feature rows.
Thinly reported entities still receive forecasts in the application, with a warning.
